\chapter{新形式的 Langlands 互反}

% ===== PDF p.046 =====
\textbf{上次}：
\begin{itemize}
\item 我们证明了分解
\[
S_k\cong\bigoplus_f\langle f\rangle_{\GL_2(\Af)}
\]
为不可约 $\GL_2(\Af)$-表示。
\item 这里 $f$ 遍历权 $k$、$\Gamma_0(N)$ 上、旁特征为 $\chi$ 的（归一化）新形式，$N$ 在变，Dirichlet 特征 $\chi:(\Z/N)^\times\to\C^\times$ 也在变。
\item 当然，在证明该分解的过程中，我们对每个不可约容许 $\GL_2(\Af)$-表示 $\pi$ 联系了一个 Hecke 特征系统
\[
\{\tilde a_p(\pi),\tilde b_p(\pi)\}_{p\notin S},
\]
其中 $S:=\{\,p\ \text{素数}:\pi^{\GL_2(\Z_p)}=0\,\}$。事实上，$\tilde a_p(\pi)$（相应地 $\tilde b_p(\pi)$）定义为双陪集
\[
\GL_2(\Z_p)\begin{pmatrix}p&0\\0&1\end{pmatrix}\GL_2(\Z_p)
\quad\text{相应地}\quad
\GL_2(\Z_p)\begin{pmatrix}p&0\\0&p\end{pmatrix}\GL_2(\Z_p)
\]
作用在由 $\GL_2(\Z_p)$ 固定的非零向量上的本征值。
\item 若 $\pi=\langle f\rangle_{\GL_2(\Af)}$，其中 $f\in S_k(\Gamma_0(N),\chi)$ 为归一化新形式，则 $S=\{\,N\ \text{的素因子}\,\}$，且
\[
\tilde a_p(\pi)=pa_p,\qquad \tilde b_p(\pi)=p^k\chi(p)
\]
对 $p\notin S$ 成立，其中 $f(q)=\sum_{n=1}^\infty a_nq^n$ 是 $f$ 的 Fourier 展开。
\end{itemize}

本讲我们将引入 $\ell$-进 Galois 表示的 Frobenius 的迹（它们是 Hecke 特征系统的「Galois 论」对应物），然后陈述 Langlands 纲领对 $S_k$ 的预测。

\textbf{$\ell$-进表示的 Frobenius 的迹}：
\begin{itemize}
\item 设 $\ell$ 为某个素数。
\item
\[
\overline\Q_\ell=\varinjlim_{E/\Q_\ell}E
\]
通过余极限拓扑成为拓扑域。
\item 设 $W$ 为有限维 $\overline\Q_\ell$-向量空间。
\item 设 $\sigma:G_\Q:=\Gal(\overline\Q/\Q)\to\GL(W)$ 为连续表示，也称为「$\ell$-进 Galois 表示」。
\item 回忆对每个素数 $p$ 存在嵌入
\[
G_{\Q_p}:=\Gal(\overline\Q_p/\Q_p)\hookrightarrow G_\Q,
\]
它在 $G_\Q$ 中相差共轭的意义下良定义。
\item 局部绝对 Galois 群 $G_{\Q_p}$ 处于正合列
\[
1\to I_p\to G_{\Q_p}\to\Gal(\overline\F_p/\F_p)=\langle\Frob_p^{\mathrm{geom}}\rangle^{\widehat\Z}\to1,
\]
其中 $I_p\subseteq G_{\Q_p}$ 是惯性子群，而
\[
\Frob_p^{\mathrm{arith}}:\overline\F_p\to\overline\F_p,\qquad x\mapsto x^p
\]
% ===== PDF p.047 =====
是 $p$ 处的算术 Frobenius。
\item 逆元 $\Frob_p^{\mathrm{geom}}:=(\Frob_p^{\mathrm{arith}})^{-1}$ 称为\textbf{几何 Frobenius}。
\item 设 $S$ 为有限素数集。称 $\sigma$ 在 $S$ 外\textbf{非分歧}，若 $\sigma(I_p)=1$ 对一切 $p\notin S$ 成立。
\item 若 $\sigma$ 在 $S$ 外非分歧，$n:=\dim_{\overline\Q_\ell}W$，则可以对 $\sigma$ 联系一个系统 $\{c_{1,p},\dots,c_{n,p}\}_{p\notin S}$，$c_{i,p}\in\overline\Q_\ell$，即
\[
c_{i,p}:=\Tr\bigl(\sigma(\Frob_p^{\mathrm{arith}})\mid\Lambda^iW\bigr).
\]
\item 注：$c_{i,p}$ 良定义，因为 $\sigma(I_p)=1$。
\item 注：$c_{i,p}$ 只依赖于 $p$，而不依赖于嵌入 $G_{\Q_p}\subseteq G_\Q$，因为迹在共轭下不变。\znote{原文此处第二条「注」把 $c_{i,p}$ 写成 $\Tr(\sigma(\Frob_p^{\mathrm{geom}})\mid\Lambda^iW)$，但 $c_{i,p}$ 的定义用的是算术 Frobenius，且 p.050 的页边注明确说用几何 Frobenius 是错的。故此处按算术 Frobenius 读，该注只保留「不依赖嵌入的选取」这一层意思。[来源：原文 p.047 的定义与 p.050 页边注。]}
\item 若 $W$ 半单，则 $\sigma$ 由 $\{c_{1,p},\dots,c_{n,p}\}_{p\notin S}$ 决定（对每个使 $W$ 在 $S$ 外非分歧的有限素数集 $S$），参见 \texttt{[Ser97, I-10]}。
\item 事实上，族 $\{c_{1,p}\}_{p\notin S}$ 已经足够（再加上 $W$ 半单且在 $S$ 外非分歧的假设）。
\item 确实：
\begin{itemize}
\item 考虑群代数 $\Lambda:=\overline\Q_\ell[G_{\Q,S}]$，其中 $G_{\Q,S}$ 是 $G_\Q$ 对由 $I_p$（$p\notin S$）生成的子群的闭包所取之商。
\item 设 $W,W'$ 为两个半单连续 $G_{\Q,S}$-表示，且
\[
\Tr(\Frob_p^{\mathrm{arith}}|W)=\Tr(\Frob_p^{\mathrm{arith}}|W')
\]
对一切 $p\notin S$ 成立。
\item 由 Chebotarev 密度定理与连续性，这蕴涵
\[
\Tr(\lambda|W)=\Tr(\lambda|W')
\]
对一切 $\lambda\in\Lambda$ 成立。
\item 对两两不同构的不可约 $\Lambda$-模 $W_1,\dots,W_m$ 与每个 $i=1,\dots,m$，存在 $\mu_i\in\Lambda$，使 $\mu_i=1$ 在 $W_i$ 上、而 $\mu_i=0$ 在 $W_j$ 上（$j\neq i$）（这是中国剩余定理的一个版本）。
\item 把
\[
W=\bigoplus_{i=1}^mW_i^{\oplus n_i},\qquad
W'=\bigoplus_{i=1}^mW_i^{\oplus n'_i}
\]
写成同型分量，其中 $W_1,\dots,W_m$ 不可约、两两不同构，$n_i,n'_i\in\N$（可能为零）。
\item 则
\[
n_i\dim_{\overline\Q_\ell}W_i=\Tr(\mu_i|W)=\Tr(\mu_i|W')=n'_i\dim_{\overline\Q_\ell}W_i
\]
对一切 $i=1,\dots,m$ 成立，即 $n_i=n'_i$，故 $W\cong W'$，如所欲。
\end{itemize}
\item 这个陈述可以看作上次提到的 $\GL_n$ 尖点自守表示「强重数一」的类比。
% ===== PDF p.048 =====
\item 重要的是我们考虑的是特征 $0$ 系数的表示；也就是说，「仅凭迹 $c_{1,p}$ 就决定（半单）表示」这一说法是错误的。\znoteplain{原文此处写作 $c_{i,p}$；但按下一句（完整的系统 $\{c_{1,p},\dots,c_{n,p}\}$ 才唯一决定）与上下文，这里指的应是单个迹 $c_{1,p}$，故译作 $c_{1,p}$。}由 Brauer--Nesbitt 定理，完整的系统 $\{c_{1,p},\dots,c_{n,p}\}_{p\notin S}$ 仍然唯一决定有限维半单表示。\znoteplain{特征 $0$ 时半单表示由迹函数唯一决定，这就是 Brauer--Nesbitt 定理的通常形式；正特征时只知道 $\Tr$ 不够，必须知道全部外幂 $\Lambda^iW$ 上的迹，这正是上面要把 $\{c_{1,p},\dots,c_{n,p}\}$ 整族留下来的原因。}
\end{itemize}

Frobenius 元的迹含有高度的算术信息。下面的讨论对理解 Langlands 纲领的算术蕴涵非常重要。

\textbf{Frobenius 的迹的算术意义}：
\begin{itemize}
\item 设 $\sigma:G_\Q\to\GL(W)$ 为 $\ell$-进表示，在 $S$ 外非分歧。
\item 迹
\[
c_{1,p}:=\Tr\bigl(\sigma(\Frob_p^{\mathrm{arith}})\mid W\bigr)
\]
编码重要的算术信息。
\item 例如，设 $F/\Q$ 为有限 Galois 扩张，Galois 群 $\Gal(F/\Q)\cong S_3$，且
\[
\sigma:G_\Q\twoheadrightarrow\Gal(F/\Q)\to\GL_2(\overline\Q_\ell)
\]
是由 $S_3$ 的唯一二维不可约表示 inflation 得到的。令 $S:=\{\,p\ \text{在}\ F\ \text{中分歧}\,\}$，并对 $p\notin S$ 记 $\Frob_{F,p}^{\mathrm{arith}}$ 为 $\Frob_p^{\mathrm{arith}}\in G_{\Q,S}$ 在 $\Gal(F/\Q)\cong S_3$ 中的像。
\item 则：
\begin{itemize}
\item $\Tr(\Frob_p^{\mathrm{arith}}|W)=2$ $\Leftrightarrow$ $\Frob_{F,p}^{\mathrm{arith}}\in\{1\}$ $\Leftrightarrow$ $p$ 在 $\OO_F$ 中完全分裂；
\item $\Tr(\Frob_p^{\mathrm{arith}}|W)=0$ $\Leftrightarrow$ $\Frob_{F,p}^{\mathrm{arith}}\in\{(1\,2),(1\,3),(2\,3)\}$ $\Leftrightarrow$ $p$ 在 $\OO_F$ 中分裂为三个不同素数；
\item $\Tr(\Frob_p^{\mathrm{arith}}|W)=-1$ $\Leftrightarrow$ $\Frob_{F,p}^{\mathrm{arith}}\in\{(1\,2\,3),(1\,3\,2)\}$ $\Leftrightarrow$ $p$ 在 $\OO_F$ 中分裂为两个不同素数。
\end{itemize}
\item 因此：知道一切（有限像的）$\ell$-进 Galois 表示的 Frobenius 的迹，就蕴涵知道一切有限扩张中非分歧素数的分解。
\item $\ell$-进表示的其它例子是平展上同调群
\[
H^i_{\text{ét}}(X_{\overline\Q},\overline\Q_\ell)
\]
对 $\Q$ 上 proper、光滑的概形 $X$ 与 $i\ge0$。
\item $G_\Q$-作用由 $H^i_{\text{ét}}(-,\overline\Q_\ell)$ 的函子性，从 $G_\Q$ 在
\[
X_{\overline\Q}=X\times_{\Spec(\Q)}\Spec(\overline\Q)
\]
上的（右）作用诱导。
% ===== PDF p.049 =====
\item proper 光滑基变换与 Grothendieck--Lefschetz 迹公式蕴涵：若 $X$ 在 $p$ 处有好约化且 $p\neq\ell$，则 $H^i_{\text{ét}}(X_{\overline\Q},\overline\Q_\ell)$ 在 $p$ 处非分歧（$i\ge0$），且
\[
\sum_{i=0}^{2\dim(X)}(-1)^i\Tr\bigl(\Frob_p^{\mathrm{geom}}\mid H^i_{\text{ét}}(X_{\overline\Q},\overline\Q_\ell)\bigr)=\sharp X(\F_p),
\]
其中 $X\to\Spec(\Z_{(p)})$ 是 $X$ 在 $p$ 处的 proper 光滑模型。
\end{itemize}

从算术的观点看，现在可以陈述 Langlands 纲领的一个要点。

\textbf{Langlands 互反\idx{Langlands 互反}的全部要点}：
\begin{itemize}
\item Frobenius 的迹应当与 Hecke 特征值相匹配！
\item 于是，算术信息被期望编码在自守信息中，而后者可能更易把握。
\item 下例取自 Matthew Emerton 的一篇帖子。他更有洞见的帖子可在其主页找到：\url{http://www.math.uchicago.edu/~emerton/}。
\begin{itemize}
\item 设 $F$ 为 $x^3-x-1$ 的分裂域。
\item 则 $F$ 是 $\Q(\sqrt{-23})$ 的 Hilbert 类域，且 $F/\Q$ 是 Galois 扩张，Galois 群为 $S_3$。
\item 像上面一样考虑相伴的二维不可约 Galois 表示 $\sigma:G_\Q\to\Gal(F/\Q)\hookrightarrow\GL_2(\overline\Q_\ell)$。
\item Hecke 证明了存在归一化新形式 $f(q)=\sum_{n=1}^\infty a_nq^n$，使
\[
a_p=\Tr\bigl(\sigma(\Frob_p^{\mathrm{arith}})\bigr)
\]
对一切满足 $p\nmid23$ 的素数 $p$ 成立。
\item 可以把 $f$ 写得更明确：考虑分歧，$f$ 必落在 $S_1(\Gamma_0(23),\chi)$ 中，其中 $\chi:(\Z/23)^\times\to\C^\times$ 是决定二次子域 $\Q(\sqrt{-23})\subseteq\Q(\zeta_{23})$ 的二次特征。
\item 熟悉模形式的人（或使用 LMFDB 之类数据库的人）会告诉我们必有
\[
f(q)=q\prod_{i=1}^\infty(1-q^i)(1-q^{23i}).
\]
\item 现在可以计算 $f$ 的展开（或在 LMFDB 中查它）：
\[
f(q)=q-q^2-q^3+\ldots+q^{58}+2q^{59}+\ldots
\]
并推断素数在 $F$ 中如何分解，例如 $59$ 完全分裂！
\end{itemize}
\item 也许 Frobenius 的迹的模性最著名的应用是 Fermat 大定理的证明。
\item 即：Fermat 定理由 Wiles/Taylor 用下述策略证明（策略由 Frey 提出，Serre、Ribet 补全）：
\begin{itemize}
\item 设 $u^p+v^p+w^p=0$，$u,v,w\in\Q$，$uvw\neq0$，$p\ge3$。
\item 考虑 $\Q$ 上由（仿射）Weierstraß 方程
\[
y^2=x(x+u^p)(x-v^p)
\]
% ===== PDF p.050 =====
（在把 $u,v,w$ 稍作调整之后）给出的椭圆曲线\idx{椭圆曲线} $E$。
\item 考虑 Galois 表示
\[
G_\Q\to\GL\bigl(H^1_{\text{ét}}(E_{\overline\Q},\overline\Q_\ell)\bigr)\cong\GL_2(\overline\Q_\ell)
\]
的对偶 $\sigma$。
\item 证明 $\sigma$ 是模的，即存在权 $2$ 的归一化新形式 $f_1(q)=\sum_{n=1}^\infty a_nq^n$，使
\[
\Tr\bigl(\sigma(\Frob_p^{\mathrm{arith}})\bigr)=a_p
\]
对几乎所有素数 $p$ 成立。
\item 用 Ribet 的一个定理可推出 $f_1$ 必与
\[
S_2(\Gamma_0(2))
\]
中的某个归一化新形式 $f_2$ 同余，即 $f_1,f_2$ 的 Fourier 系数是代数整数，且在某个素数模下同余。
\item 但 $S_2(\Gamma_0(2))=0$，故 $f_2$ 不存在。
\item 于是得到所要的矛盾。
\item 更多细节参见 \texttt{[Wil95]}、\texttt{[Rib90]}。
\end{itemize}
\end{itemize}

现在陈述关于 Galois 表示与尖点形式之间关系的预期。

\textbf{新形式的 Langlands 互反}：
\begin{itemize}
\item 固定素数 $\ell$ 与同构
\[
\iota:\overline\Q_\ell\xrightarrow{\ \sim\ }\C.
\]
\item 对新形式，Langlands 纲领结合 Fontaine--Mazur 猜想预测存在双射
\[
\mathrm{LL}:\cA_{\mathrm{mod}}\xrightarrow{\ \sim\ }\mathcal G_{\mathrm{mod}},
\]
从集合
\[
\cA_{\mathrm{mod}}:=\Bigl\{\ \text{不可约}\ \GL_2(\Af)\text{-子表示}\ \pi\subseteq\bigoplus_{k\ge1}S_k\ \Bigr\}
\]
到某个集合 $\mathcal G_{\mathrm{mod}}$，后者由不可约二维 Galois 表示
\[
\sigma:G_\Q\to\GL_2(\overline\Q_\ell)
\]
组成。
\item $\mathrm{LL}$ 应满足：对一切 $\pi\in\cA_{\mathrm{mod}}$，令 $\sigma:=\mathrm{LL}(\pi)$，有
\[
\iota\bigl(\Tr(\sigma(\Frob_p^{\mathrm{arith}}))\bigr)=a_p(\pi),\qquad
\iota\bigl(\det(\sigma(\Frob_p^{\mathrm{arith}}))\bigr)=b_p(\pi)
\]
对某个指定的有限素数集 $S$ 之外的 $p$ 成立。\znote{原文此处页边注「In a previous version I considered the geometric Frobenius, which was wrong」重复排印了两次（同一句话上下各一遍），译文只保留一次。这句话针对的正是上面的匹配公式：在二维情形 $\Frob_p^{\mathrm{geom}}=(\Frob_p^{\mathrm{arith}})^{-1}$，而 $\Tr(A^{-1})=\Tr(A)/\det(A)$，若误用几何 Frobenius，得到的将是 $a_p/b_p$ 而不是 $a_p$。[来源：原文 p.050 页面的页边注与 p.047 的 Frobenius 定义。]}
\item 这里使用的是本征值 $a_p(\pi)$（相应地 $b_p(\pi)$），而不是 $\tilde a_p(\pi)$、$\tilde b_p(\pi)$。\znoteplain{固定 $\iota:\overline\Q_\ell\cong\C$ 的作用是把 $\ell$-进数（Frobenius 的迹与行列式）搬到复数域，才能与 Fourier 系数 $a_p(\pi),b_p(\pi)\in\C$ 比较。换一个 $\iota$ 会把所有这些代数数同时换成它们的共轭，双射 $\mathrm{LL}$ 也随之改变，所以 $\iota$ 必须一开始就固定。}
\item 注：这样的双射 $\mathrm{LL}$ 若存在，则由相应的重数一定理唯一决定。
\item 猜想地，集合 $S$ 有明确的描述。回忆素数 $p$ 对 $\sigma$（相应地 $\pi$）非分歧，若
\[
\sigma(I_p)=1
\qquad\text{相应地}\qquad
\pi^{\GL_2(\Z_p)}\neq0.
\]
% ===== PDF p.051 =====
\item 则猜想地，$p\neq\ell$ 对 $\pi$ 非分歧当且仅当它对 $\sigma:=\mathrm{LL}(\pi)$ 非分歧。
\item Hecke 特征值的匹配更精确地是说：$\sigma(\Frob_p^{\mathrm{arith}})$ 的特征多项式应为
\[
X^2-\iota^{-1}(a_p(\pi))X+\iota^{-1}(b_p(\pi)),
\]
对 $\pi$ 的每个非分歧素数 $p\neq\ell$。
\item 课程的下一目标是说明 $\mathrm{LL}$ 可以如何构造，即如何把 Galois 表示联系到归一化新形式。为此我们遵循 Deligne 的结果（参见 \texttt{[Del71b]}）与 Deligne/Serre 的结果（参见 \texttt{[DS74]}）。
\end{itemize}

下面更精确地说明预期与一般新形式相伴的是哪个 Galois 表示集合。

\textbf{集合 $\mathcal G_{\mathrm{mod}}$}：
\begin{itemize}
\item 集合 $\mathcal G_{\mathrm{mod}}$ 由 Fontaine/Mazur 在他们关于几何 Galois 表示的著名猜想中定义，参见 \texttt{[FM95]}。
\item 设 $\sigma:G_\Q\to\GL_2(\overline\Q_\ell)$ 为不可约二维 $\ell$-进表示。
\item 则按定义 $\sigma\in\mathcal G_{\mathrm{mod}}$ 当且仅当
\begin{itemize}
\item $\sigma$ 在几乎所有 $p$ 处非分歧，即 $\sigma(I_p)=1$ 对某个有限素数集 $S$ 之外的 $p$ 成立；
\item $\sigma$ 是\textbf{奇的}（odd），即对每个复共轭 $c\in\Gal(\C/\R)\subseteq G_\Q$ 有
\[
\det(\sigma(c))=-1;
\]
\item 限制
\[
\sigma_\ell:=\sigma|_{G_{\Q_\ell}}:G_{\Q_\ell}\to\GL_2(\overline\Q_\ell)
\]
是 de Rham 的，Hodge--Tate 权\idx{Hodge--Tate 权}为 $0,\omega$（某个 $\omega\in\N$；注意 $\omega=0$ 是允许的）。
\end{itemize}
\item 若 $\sigma=\mathrm{LL}(\pi)$ 且 $\pi\subseteq S_k$，则猜想地 $\omega=k-1$。
\item 利用许多人的工作（Wiles、Kisin、Emerton、Skinner--Wiles、Pan、$\dots$），任何 Hodge--Tate 权互不相同的 $\sigma\in\mathcal G_{\mathrm{mod}}$ 在 $\ell\ge5$ 时是模的，参见 \texttt{[Pan19, Theorem 1.0.4.]}。
\item 若 $\sigma\in\mathcal G_{\mathrm{mod}}$ 有有限像，则 $\sigma$ 是模的，参见 \texttt{[PS16]}、\texttt{[GH19, Section 13.4.]}（以及其中的参考文献）。
\item 猜想地：全部 Hodge--Tate 权为 $0$ $\Leftrightarrow$ $\sigma$ 的像有限。
\end{itemize}

我们必须对 de Rham 这一重要概念作一段简短的离题讨论。

\textbf{de Rham 表示\idx{de Rham 表示}}：
\begin{itemize}
\item 设 $p$ 为素数（即上面的 $\ell$）。
\item 设 $K/\Q_p$ 为离散赋值、非 Archimedes 扩张，剩余域完美（例如 $K/\Q_p$ 有限）。
\item 设 $\sigma:G_K:=\Gal(\overline K/K)\to\GL_m(\Q_p)$ 为连续表示。
\item 则 $\sigma$ 的系数落在某个有限扩张 $E/\Q_p$ 中（因为 $G_K$ 紧，而这里 $\Q_p$ 带余极限拓扑）。
\item 于是得到 $\Q_p$-线性表示，即「局部 $p$-进 Galois 表示」
\[
\rho:G_K\to\GL_m(E)\subseteq\GL_n(\Q_p),\qquad n:=m\cdot\dim_{\Q_p}E.
\]
% ===== PDF p.052 =====
\item 于是可以改变设置，令 $V$ 为有限维 $\Q_p$-向量空间，$\rho:G_K\to\GL(V)$ 为连续表示。
\item 则 $V$ 称为 \textbf{de Rham 的}，若
\[
\dim_K\bigl(B_{\mathrm{dR}}\otimes_{\Q_p}V\bigr)^{G_K}=\dim_{\Q_p}V
\]
对某个带 $G_K$-作用的域扩张 $B_{\mathrm{dR}}$ 成立，参见 \texttt{[BC09, Section 6]}。
\item 即：$B_{\mathrm{dR}}$ 是 Fontaine 的 $p$-进周期域，参见 \texttt{[Fon94]}、\texttt{[BC09, Definition 4.4.7]}。抽象地，
\[
B_{\mathrm{dR}}\cong C_K((t)),
\]
其中 $C_K$ 是 $K$ 对其 $p$-进赋值的完备化（这个拓扑比 $K$ 上的余极限拓扑更粗）。
\item $B_{\mathrm{dR}}$ 是离散赋值域，剩余域为 $C_K$；由此得到的递减滤过 $\mathrm{Fil}^\bullet B_{\mathrm{dR}}$ 是 $G_K$-稳定的，且作为 $C_K$-半线性 $G_K$-表示
\[
\mathrm{gr}^\bullet B_{\mathrm{dR}}\cong\bigoplus_{j\in\Z}C_K(j),
\]
其中 $C_K(j):=C_K\otimes_{\Z_p}\Z_p(j)$，而
\[
\Z_p(1):=T_p\mu_{p^\infty}(C_K)
\]
是与分圆特征 $\chi_{\mathrm{cyc}}:G_K\to\Z_p^\times\cong\Aut(\mu_{p^\infty}(C_K))$ 相伴的 $p$-进 Galois 表示。
\item Tate 的一个重要定理（参见 \texttt{[BC09, Theorem 2.2.7]}）说：设 $\chi:G_K\to\Q_p^\times$ 为特征。则
\[
H^0\bigl(G_K,C_K(\chi)\bigr)=
\begin{cases}
K,&\chi|_{I_K}\ \text{的像有限},\\
0,&\text{否则},
\end{cases}
\]
这里 $I_K\subseteq G_K$ 是惯性子群。
\item 特别地：
\[
H^0\bigl(G_K,C_K(j)\bigr)=0,\qquad j\neq0,
\]
因为 $j\neq0$ 时 $\chi_{\mathrm{cyc}}^j|_{I_K}$ 的像无限。
\item 结论：若 $\dim_{\Q_p}V=1$，则 $V$ 是 de Rham 的当且仅当 $V\cong\chi\cdot\chi_{\mathrm{cyc}}^j$，其中 $\chi:G_K\to\Q_p^\times$ 是某个使 $I_K$ 的像有限的特征，$j\in\Z$。特别地，存在许多非 de Rham 的表示，例如特征 $(\chi_{\mathrm{cyc}}^{p-1})^a$，$a\in\Z_p\setminus\Z$。\znote{原文作 $\dim_\C V=1$；这里 $V$ 是 $G_K$ 的 $\Q_p$-表示，而 de Rham 判据比较的是 $\dim_K(B_{\mathrm{dR}}\otimes_{\Q_p}V)^{G_K}$ 与 $\dim_{\Q_p}V$，故应读作 $\dim_{\Q_p}V=1$，否则整个陈述没有意义。[来源：原文 p.052 上一条 de Rham 定义。]}
\item 我们提一些关于 de Rham 表示的一般事实。
\begin{itemize}
\item 若存在 $K$ 的有限扩张 $K'$ 使 $\rho|_{G_{K'}}$ 是 de Rham 的，则 $\rho:G_K\to\GL(V)$ 是 de Rham 的，参见 \texttt{[BC09, Proposition 6.3.8]}。
\item de Rham 表示对子商、对偶与张量积稳定，参见 \texttt{[BC09, Section 6.1]}。
\end{itemize}
\item 设 $V$ 为 $n$ 维 de Rham 表示。则
\[
V\otimes_{\Q_p}C_K\cong C_K(j_1)\oplus\ldots\oplus C_K(j_n)
\]
作为 $C_K$-半线性 $G_K$-表示，其中 $j_1,\dots,j_n\in\Z$。无序族 $(j_1,\dots,j_n)$ 称为 $V$ 的 \textbf{Hodge--Tate 权}之族。
\end{itemize}

% ===== PDF p.053 =====
因此，de Rham 条件给 $p$-进表示加上了某种「整性」条件。

我们以说明 $\mathrm{LL}$ 如何构造（当 $k=2$ 时）来结束本讲。

\textbf{构造 $\mathrm{LL}$ 的策略（粗略）}：
\begin{itemize}
\item 设 $\pi\in\cA_{\mathrm{mod}}$，相伴的 Hecke 特征系统为 $\{a_p,b_p\}_{p\notin S}$。
\item 为简单起见设 $\pi\subseteq S_2$，即 $\pi$ 由某个权 $2$ 的新形式生成。
\item 则我们将构造（接近）一个 $\GL_2(\Af)$-等变嵌入
\[
S_2\hookrightarrow H^1\bigl(\GL_2(\Q)\backslash(\GL_2(\Af)\times\HH^\pm),\C\bigr)
:=\varinjlim_{K}H^1\bigl(\GL_2(\Q)\backslash(\GL_2(\Af)/K\times\HH^\pm),\C\bigr),
\]
把 $S_2$ 嵌入上同调。
\item 给定的同构 $\iota:\overline\Q_\ell\cong\C$ 给出上同调群的同构
\[
H^1\bigl(\GL_2(\Q)\backslash(\GL_2(\Af)\times\HH^\pm),\C\bigr)
\cong H^1\bigl(\GL_2(\Q)\backslash(\GL_2(\Af)\times\HH^\pm),\overline\Q_\ell\bigr).
\]
\item 如前所述，对（充分小的）紧开子群 $K\subseteq\GL_2(\Af)$，复流形
\[
\GL_2(\Q)\backslash(\GL_2(\Af)/K\times\HH^\pm)
\]
实际上是代数的，即由某个拟射影概形 $\widetilde X_K\to\Spec(\C)$ 的 $\widetilde X_K(\C)$ 给出。
\item 平展比较定理蕴涵
\[
H^1\bigl(\GL_2(\Q)\backslash(\GL_2(\Af)/K\times\HH^\pm),\overline\Q_\ell\bigr)
\cong H^1_{\text{ét}}(\widetilde X_K,\overline\Q_\ell).
\]
\item 现在奇迹发生了：对（充分小的）紧开子群 $K\subseteq\GL_2(\Af)$，拟射影概形 $\widetilde X_K\to\Spec(\C)$ 典范地定义在 $\Spec(\Q)$ 上！换言之，$\widetilde X_K\cong X_K\times_{\Spec(\Q)}\Spec(\C)$，其中 $X_K\to\Spec(\Q)$ 是（相容的）拟射影概形。
\item 特别地，由平展上同调在代数闭基域变换下的不变性，
\[
H^1_{\text{ét}}(\widetilde X_K,\overline\Q_\ell)\cong H^1_{\text{ét}}(X_{K,\overline\Q},\overline\Q_\ell).
\]
但现在右端带有 $G_\Q$ 的作用！
\item 现在看 $G_\Q\times\GL_2(\Af)$-模
\[
H^1_{\text{ét}}(X_{\overline\Q},\overline\Q_\ell)
:=\varinjlim_{K}H^1_{\text{ét}}(X_{K,\overline\Q},\overline\Q_\ell).
\]
\item 则我们将验证 $G_\Q$-模
\[
\mathrm{LL}(\pi):=\Hom_{\GL_2(\Af)}\bigl(\pi,H^1_{\text{ét}}(X_{\overline\Q},\overline\Q_\ell)\bigr)
\]
是二维的，并且可以用 Hecke 特征值表达 Frobenius 的迹。这个二维表示的对偶将是所要的表示。
\item 把常值层 $\overline\Q_\ell$ 换成某个局部系统，类似的策略对权 $k\ge2$ 也有效。\znote{原文 "Replacing $\overline\Q_\ell$ by some local system"：这里换掉的是系数层，即把常值层 $\overline\Q_\ell$ 换成与 $\omega^{\otimes k}$ 相应的局部系统（而不是把表示 $\sigma$ 换掉）；只有这样从 $H^1$ 中切出的才是权 $k$ 的表示。[来源：原文 p.053 与第 9 讲对 $\omega^{\otimes k}$ 的用法。]}
\end{itemize}

% ===== PDF p.054 =====
\begin{itemize}
\item 然而这对 $k=1$ 没有帮助。这里必须利用：权 $1$ 的模形式可以与某个权 $k\ge2$ 的模形式同余，然后使用前一种情形。
\end{itemize}
